\section{Extended Theoretical Statements}
\label{app:analysis-statements}

We separate exact algebraic properties from assumptions used to make the
mixed-bit objective separable.  Full proofs are in \cref{app:proofs}.

\subsection{Exactness and balanced score geometry}

\paragraph{Proposition 1 (exact biorthogonal score preservation).}
Let $\widetilde C_q,C_k\in\R^{d\times d}$ be positive definite and let
$A,D$ be defined by \cref{eq:qk-svd,eq:qk-factors}.
Then
\begin{equation}
 AD^\top=I.
 \label{eq:prop-biorthogonal}
\end{equation}
Consequently, for every $q,k\in\R^d$, not only for calibration samples,
\begin{equation}
 (A^\top q)^\top(D^\top k)=q^\top k.
 \label{eq:prop-score-exact}
\end{equation}
Thus \method's transform is not a low-rank approximation and adds no
full-precision score error.

\paragraph{Proposition 2 (simultaneous second-moment balance).}
Under the same conditions,
\begin{equation}
 A^\top\widetilde C_qA=D^\top C_kD=\Sigma.
 \label{eq:prop-cov-balance}
\end{equation}
If transformed Queries and Keys are paired independently according to these
second moments, their expected squared dot product is
\begin{equation}
 \mathbb E[({q'}^\top k')^2]=\tr(\Sigma^2)
 =\sum_{\ell=1}^{d}\sigma_\ell^2.
 \label{eq:score-energy}
\end{equation}
Therefore the singular order exposes joint score energy: retaining the first
$r$ unquantized coordinates leaves expected tail score energy
$\sum_{\ell>r}\sigma_\ell^2$ under the construction moments.
With shrinkage and the spectral floor,
\cref{eq:prop-cov-balance} holds exactly for the construction moments
$\widetilde C_q$ and $C_k$, while the empirical moments differ by
\begin{equation}
 \Delta_q^{\rm mom}
 =A^\top(\widehat C_q-\widetilde C_q)A,\qquad
 \Delta_k^{\rm mom}
 =D^\top(\widehat C_k-C_k)D.
 \label{eq:shrinkage-residual}
\end{equation}
These residuals are deterministic and measurable.  The exact score identity
in Proposition~1 is unaffected because it depends on $AD^\top=I$, not on
either empirical second moment being equal to its construction moment.

\paragraph{Theorem (optimal rank-$r$ score subspace).}
Let independent $q$ and $k$ have positive-definite uncentered second moments
$C_q$ and $C_k$, and approximate $q^\top k$ by $q^\top B_r k$ with
$\operatorname{rank}(B_r)\le r$.  If
$C_q^{1/2}C_k^{1/2}=U\Sigma V^\top$, then
\begin{equation}
 B_r^\star=C_q^{-1/2}U_r\Sigma_rV_r^\top C_k^{-1/2}
 =A_rD_r^\top
 \label{eq:optimal-rank-score-map}
\end{equation}
globally minimizes the expected squared score error, and
\begin{equation}
 \min_{\operatorname{rank}(B_r)\le r}
 \mathbb E[(q^\top k-q^\top B_rk)^2]
 =\sum_{\ell>r}\sigma_\ell^2.
 \label{eq:optimal-rank-score-error}
\end{equation}
This is an Eckart--Young result in the Query--Key metric.  It gives a precise
meaning to ``joint score energy'': before quantization, the leading
QK-balanced coordinates are the optimal rank-$r$ bilinear approximation,
whereas Key-PCA is optimal for a different Key-reconstruction objective.
The theorem assumes independent draws from the construction moments; the
finite Cartesian-sample identity below is what the allocator actually
measures.

\paragraph{Corollary (same-rate cost of stale QK coordinates).}
For condition $c$, let
$M_c=C_{q,c}^{1/2}C_{k,c}^{1/2}$ have singular values
$\{\sigma_{c,j}\}$ and let $B_{c_0}$ be any rank-$r$ map frozen under
another condition.  Then
\begin{align}
 L_c(B_{c_0})-L_c(B_c^\star)
 &=
 \|M_c-C_{q,c}^{1/2}B_{c_0}C_{k,c}^{1/2}\|_F^2
 -\sum_{j>r}\sigma_{c,j}^2
 \geq0 .
 \label{eq:stale-qk-coordinate-excess}
\end{align}
With a strict $r$th singular gap, equality requires the frozen map to
realize the current leading singular subspaces.  Thus request-local
coordinates can strictly reduce score distortion at identical rank and
rate.  Mixed-bit QKSieve is not a pure rank truncation, but its zero- and
low-bit tail inherits this mechanism.

\paragraph{Proposition (residual under paired Query--Key dependence).}
Independence is sufficient for the preceding rank result but need not hold
for Query--Key pairs observed in attention.  Define
\begin{equation}
 H=\mathbb E[(k\otimes q)(k\otimes q)^\top],\qquad
 H_0=C_k\otimes C_q,\qquad
 \Delta_H=H-H_0 .
\end{equation}
For any bilinear map $B$, let
$L_{\rm pair}(B)=\mathbb E[(q^\top(I-B)k)^2]$ under the paired law and let
$L_{\rm ind}(B)$ be the same loss under independent draws with the same
second moments.  Then
\begin{equation}
 |L_{\rm pair}(B)-L_{\rm ind}(B)|
 \le \|\Delta_H\|_{\op}\|I-B\|_F^2.
 \label{eq:paired-qk-dependence}
\end{equation}
Thus second moments alone cannot certify rank optimality for arbitrary
dependent pairs.  The bound isolates the missing fourth-order statistic;
our audit reports paired-versus-Cartesian loss residuals for the frozen
candidate maps rather than assuming it is zero.

\paragraph{Corollary (Key-PCA is the isotropic-Query special case).}
If $C_q=\alpha I$ and
$C_k=V\Lambda V^\top$, then the left and right singular vectors in
\cref{eq:qk-svd} are both $V$, with
$\Sigma=\alpha^{1/2}\Lambda^{1/2}$.
Consequently, QK-balanced coordinates have the same directions and ordering
as Key-PCA, differing only by reciprocal Query/Key scalings that preserve
dot products.  QK balancing therefore reduces to Key-PCA when all Query
directions are equally important; its additional rotation and allocation
are needed precisely when the Query moment is anisotropic.

\subsection{Why query-weighted error is the correct allocation target}

\paragraph{Proposition 3 (empirical logit-MSE identity).}
Let $\{q'_j\}_{j=1}^{m_q}$ be any Query sample and
$\{e_i\}_{i=1}^{m_k}$ any Key-quantization errors.
Define their uncentered second moments
$C'_q=m_q^{-1}\sum_jq'_j{q'_j}^\top$ and
$C_e=m_k^{-1}\sum_ie_ie_i^\top$.
Then the Cartesian empirical score MSE is exactly
\begin{equation}
 \frac{1}{m_qm_k}\sum_{j,i}({q'_j}^\top e_i)^2
 =\tr(C'_qC_e).
 \label{eq:trace-qmse}
\end{equation}
No zero-mean, Gaussian, or sample-independence assumption is needed for this
finite-sample identity.
For normalized attention logits in \cref{eq:full-attn}, the MSE is the
right-hand side divided by $d$; this common factor does not change the
allocation minimizer.

Partitioning the coordinates into bands gives
\begin{equation}
 \tr(C'_qC_e)=
 \sum_g\tr(C'_{q,gg}C_{e,gg})
 +\sum_{g\ne h}\tr(C'_{q,gh}C_{e,hg}).
 \label{eq:block-qmse}
\end{equation}
The production allocator minimizes the first sum.
This is exact when either moment is block diagonal in the QK-balanced bands.
Otherwise the omitted cross-band term obeys
\begin{equation}
 \left|\sum_{g\ne h}\tr(C'_{q,gh}C_{e,hg})\right|
 \le\sum_{g\ne h}
 \|C'_{q,gh}\|_F\|C_{e,hg}\|_F.
 \label{eq:cross-band-bound}
\end{equation}
The right-hand side is an auditable diagnostic for the separability
approximation.

\paragraph{Corollary (rate-constrained dominance).}
Let $\mathcal B_R$ be the finite set of band allocations satisfying the
physical rate budget in \cref{eq:qmse-allocation}.  Because the dynamic
program solves the separable objective globally,
\begin{equation}
 \sum_gD_g(b_g^\star)\le\sum_gD_g(\bar b_g)
 \quad\text{for every }\bar b\in\mathcal B_R.
 \label{eq:allocation-dominance}
\end{equation}
Thus automatic allocation cannot have larger calibration distortion than
any feasible fixed allocation under the same quantizers and separable
objective.  This statement does not by itself guarantee lower held-out
distortion; \cref{eq:query-drift-bound,eq:finite-query-bound} quantify the
two ways calibration dominance can fail to transfer.

\paragraph{Proposition (one versus two joint-spectrum weights).}
Assume the QK-balanced construction moments are exact, so that
$\mathbb E[q'q'^\top]=\mathbb E[k'k'^\top]
=\Sigma=\operatorname{diag}(\sigma_1,\ldots,\sigma_d)$.
For an idealized coordinatewise quantizer, assume zero-mean uncorrelated
errors with
$\mathbb E[e_j^2]=c_j(b_j)\sigma_j$, where $c_j(b_j)\ge0$ is its
rate-dependent relative distortion.  Then its Key-MSE is
\begin{equation}
 J_{\rm K}(b)=\tr(C_e)
 =\sum_{j=1}^{d}c_j(b_j)\sigma_j,
 \label{eq:keymse-joint-spectrum}
\end{equation}
whereas its calibration logit-MSE is
\begin{equation}
 J_{\rm QK}(b)=\tr(\Sigma C_e)
 =\sum_{j=1}^{d}c_j(b_j)\sigma_j^2.
 \label{eq:qmse-joint-spectrum}
\end{equation}
Thus Key-MSE in QK-balanced coordinates already carries one factor of the
joint Query--Key spectrum; it is not equivalent to Key-MSE in a
Query-independent Key-PCA basis.  Logit-MSE applies a second factor and can
concentrate rate more strongly on the leading spectrum.  The proposition
does not rank the two objectives for held-out decoding: block quantization,
nonproportional error, Query drift, and ranking margins can favor either.

\paragraph{Theorem (held-out regret under Query drift).}
For allocation $b$, let
\begin{align}
 \widetilde J_t(b)
 &=\sum_g\tr(C'_{q,t,gg}C_{e,gg}(b_g)),\\
 \Gamma_t(b)
 &=\sum_{g\ne h}\tr(C'_{q,t,gh}C_{e,hg}(b)),
 \qquad J_t(b)=\widetilde J_t(b)+\Gamma_t(b),
\end{align}
where $t\in\{\mathrm{cal},\mathrm{dec}\}$.
Define $\Delta=C'_{q,\mathrm{dec}}-C'_{q,\mathrm{cal}}$ and
\begin{equation}
 \Omega(b)=\sum_g\|\Delta_{gg}\|_{\op}
                    \tr(C_{e,gg}(b_g)).
 \label{eq:allocation-drift-penalty}
\end{equation}
For the calibration minimizer $b^\star$ and every fixed feasible allocation
$\bar b$,
\begin{equation}
 J_{\mathrm{dec}}(b^\star)-J_{\mathrm{dec}}(\bar b)
 \le \Omega(b^\star)+\Omega(\bar b)
 +|\Gamma_{\mathrm{dec}}(b^\star)|
 +|\Gamma_{\mathrm{dec}}(\bar b)|.
 \label{eq:heldout-allocation-regret}
\end{equation}
Thus the automatic allocation is near-optimal on held-out decode Queries
when diagonal-block covariance drift and cross-band error interactions are
small.  Every term on the right is measurable without task labels.  The
result also states precisely why calibration optimality alone is
insufficient.

\paragraph{Idealized rate-distortion interpretation.}
Suppose a continuous high-rate surrogate has band distortion
$D_g(b_g)=\alpha_g2^{-2b_g}$, ignores metadata, and permits $b_g\ge0$.
The exact solution of
$\min\sum_gD_g(b_g)$ subject to $\sum_gb_g\le R$ is
\begin{equation}
 b_g^\star=\left[\frac12\log_2\frac{\alpha_g}{\tau}\right]_+,
 \qquad
 \sum_g b_g^\star=R,
 \label{eq:waterfilling-bits}
\end{equation}
for a threshold $\tau$ fixed by the rate constraint.
When every band is active,
$b_g^\star=R/G+\frac12\log_2(\alpha_g/\bar\alpha)$, where
$\bar\alpha$ is the geometric mean of the $\alpha_g$.
The formula predicts more bits for bands with larger Query-weighted score
sensitivity.  It is an interpretation, not the production optimizer:
\method's discrete program evaluates the actual 0/1/2/4/8-bit quantizers and
scale metadata exactly.

\paragraph{Corollary (heterogeneity gap against uniform rate).}
Under the same all-active high-rate model, compare the optimum above with
the uniform allocation $b_g=R/G$.  Their distortions satisfy
\begin{equation}
 \frac{J_{\mathrm{uniform}}}{J_{\mathrm{mixed}}^\star}
 =
 \frac{\frac1G\sum_{g=1}^{G}\alpha_g}
      {(\prod_{g=1}^{G}\alpha_g)^{1/G}}
 \ge 1,
 \label{eq:uniform-mixed-ratio}
\end{equation}
with equality if and only if every band has identical score sensitivity.
Thus the idealized gain from mixed rate is exactly the
arithmetic-to-geometric-mean ratio of Query-weighted band sensitivities.
This yields a falsifiable prediction: layer--head groups with larger
sensitivity heterogeneity should show larger uniform-to-automatic QK-MSE
gains.  The statement compares rates in a fixed coordinate system; our
uniform-bit QK-balanced ablation isolates this effect, while the FIER
comparison additionally changes the coordinate system.

\subsection{Conditional fixed-rate allocation and length}

Let a condition
$c=(x,N,\ell,h,t)$ specify request, history length, layer, KV head, and
decode position.  A policy $\theta\in\Theta_R$ contains a score-preserving
QK transform and a feasible mixed-bit allocation with physical index rate
at most $R$; assume the evaluated policy class contains a minimizer.  For
exact and proxy scores, write
$\xi_i(\theta)=\widehat s_i(\theta)-s_i$,
$\bar\xi=N^{-1}\sum_i\xi_i$, and
\begin{equation}
 \eta_c^2(\theta)=
 \frac1N\sum_{i=1}^{N}
 [\xi_i(\theta)-\bar\xi(\theta)]^2 .
\end{equation}
Let $m_{r,B}(c)=s_{(r)}-s_{(B+1)}>0$ for $r\le B$.

\paragraph{Theorem (conditional fixed-rate core certificate).}
Define
\begin{equation}
 \mathcal C_R(c)=
 \min_{\theta\in\Theta_R}
 \frac{4N\eta_c^2(\theta)}{m_{r,B}^2(c)} .
 \label{eq:fixed-rate-core-certificate}
\end{equation}
If $\mathcal C_R(c)<1$, there exists a policy of rate at most $R$ whose
proxy top-$B$ contains exact top-$r$.  More generally, the minimizing policy
misses at most
\begin{equation}
 \min\{r,\lfloor\mathcal C_R(c)\rfloor\}
 \label{eq:conditional-fixed-rate-misses}
\end{equation}
exact top-$r$ tokens.  Consequently, if
$\mathcal C_R(c)<1$ for every $c$ in a condition family, one common rate
$R$ suffices for the family even when the minimizing
$\theta_c^\star$ differs across requests, lengths, layers, or heads.
The statement does not imply that a single frozen policy, or rate $R=240$,
is sufficient for arbitrary conditions.

\paragraph{Theorem (conditional policies dominate frozen templates).}
Let $L(c,\theta)\geq0$ be any integrable condition-specific loss and use
the same rate-constrained policy class $\Theta_R$ on both sides.  Then
\begin{equation}
 \underbrace{\mathbb E_c\!
  \left[\min_{\theta\in\Theta_R}L(c,\theta)\right]}
 _{\mathcal L_{\rm adapt}(R)}
 \le
 \underbrace{\min_{\theta\in\Theta_R}
  \mathbb E_c[L(c,\theta)]}
 _{\mathcal L_{\rm frozen}(R)} .
 \label{eq:conditional-policy-dominance}
\end{equation}
For finite $\Theta_R$, equality holds if and only if an optimal frozen
policy is conditionwise optimal almost surely.  Otherwise the inequality
is strict.  Thus adaptation can reduce the rate required for a target
risk without increasing the per-condition rate; it cannot make an
insufficient feasible policy class sufficient.

In practice QKSieve minimizes an estimated loss
$\widehat L(c,\theta)$.  If
$\widehat\theta(c)\in\arg\min_\theta\widehat L(c,\theta)$, then
\begin{equation}
 L(c,\widehat\theta(c))-\min_{\theta\in\Theta_R}L(c,\theta)
 \le
 2\sup_{\theta\in\Theta_R}
 |\widehat L(c,\theta)-L(c,\theta)|.
 \label{eq:conditional-selection-regret}
\end{equation}
This separates the value of conditional allocation from finite-sample
calibration and Query-distribution drift.

\paragraph{Theorem (probabilistic crossing bound).}
Condition on the exact scores.  Suppose the centered per-token score errors
are independent, zero-mean, and sub-Gaussian with variance proxy at most
$\sigma_c^2(\theta)$.  Then
\begin{equation}
 \Pr\!\left[
  T_r(c)\not\subseteq\widehat T_B(\theta,c)
 \right]
 \le
 r(N-B)
 \exp\!\left[
  -\frac{m_{r,B}^2(c)}{4\sigma_c^2(\theta)}
 \right].
 \label{eq:crossing-probability}
\end{equation}
Thus failure probability at most $\delta$ is certified by
\begin{equation}
 \sigma_c^2(\theta)
 \le
 \frac{m_{r,B}^2(c)}
 {4\log(r(N-B)/\delta)} .
 \label{eq:crossing-subgaussian-condition}
\end{equation}
This average-case result complements the assumption-free deterministic
certificate above; quantization errors in a realized request are not claimed
to be independent.

\paragraph{A separable crossing-risk objective.}
For $i\in T_r$, $j\notin T_B$, let
$\Delta_{ij}=s_i-s_j>0$, and decompose score error by band as
$\xi_i=\sum_g\xi_{i,g}$.  Markov's inequality and a union bound give
\begin{equation}
 \Pr[T_r\not\subseteq\widehat T_B]
 \le
 \sum_{\substack{i\in T_r\\j\notin T_B}}
 \frac{\mathbb E[(\xi_j-\xi_i)^2]}{\Delta_{ij}^2}.
 \label{eq:crossing-markov-objective}
\end{equation}
When cross-band error terms vanish, the right-hand side decomposes as
\begin{align}
 J_{\rm cross,c}(b)
 &=\sum_g D_{{\rm cross},c,g}(b_g),\\
 D_{{\rm cross},c,g}(b)
 &=\sum_{\substack{i\in T_r\\j\notin T_B}}
 \frac{\mathbb E[(\xi_{j,g}(b)-\xi_{i,g}(b))^2]}
 {\Delta_{ij}^2}.
 \label{eq:crossing-band-distortion}
\end{align}
The same finite dynamic program as \cref{eq:qmse-allocation} can therefore
minimize a margin-weighted crossing surrogate under the unchanged physical
rate.  Nonzero cross-band terms admit the same Frobenius audit as
\cref{eq:cross-band-bound}.

\paragraph{Corollary (condition-dependent water filling).}
Under the continuous all-active surrogate
$D_{{\rm cross},c,g}(b_g)=\alpha_{c,g}2^{-2b_g}$ and
$\sum_gb_g=R$,
\begin{equation}
 b_{c,g}^\star=
 \frac RG+
 \frac12\log_2
 \frac{\alpha_{c,g}}
 {(\prod_u\alpha_{c,u})^{1/G}} .
 \label{eq:conditional-waterfilling}
\end{equation}
The total rate is unchanged, but the optimal band allocation changes
whenever the relative crossing sensitivities $\alpha_{c,g}$ change.  A
uniform change in margin tightens the achievable certificate but does not,
by itself, change relative band rates; request- or length-induced
nonuniform sensitivity is what moves bits between bands.

\paragraph{Corollary (global layer--head--band allocation).}
Index units by $a=(\ell,h,g)$, let $w_a$ be physical rate per scalar bit,
and consider
\begin{equation}
 \min_{b_a\geq0}\sum_a\alpha_{c,a}2^{-2b_a}
 \quad\text{s.t.}\quad
 \sum_aw_ab_a\leq R_{\rm total}.
\end{equation}
The continuous optimum is
\begin{equation}
 b_{c,a}^{\star}
 =
 \left[
 \frac12\log_2\frac{\alpha_{c,a}}{\tau_cw_a}
 \right]_+ ,
 \label{eq:global-conditional-waterfilling}
\end{equation}
where $\tau_c$ enforces the unchanged total rate.  Separate per-head rate
constraints give one multiplier per $(\ell,h)$; a ragged global packed
layout permits one multiplier across all units.

\paragraph{Output-aware source of conditional sensitivity.}
For one head, $o=\sum_ip_iv_i$ and $p=\softmax(s)$.  The exact first-order
term under score perturbation is
\begin{equation}
 \delta o
 =
 \sum_i p_i(v_i-o)\delta s_i
 +O(\|\delta s\|_2^2).
 \label{eq:attention-local-jacobian}
\end{equation}
Consequently, the bandwise local output distortion is
\begin{equation}
 D_{{\rm out},c,\ell,h,g}(b)
 =
 \mathbb E\left\|
  \sum_i p_i(v_i-o)\delta s_{i,g}(b)
 \right\|_2^2 .
 \label{eq:output-aware-band-distortion}
\end{equation}
If tokenwise errors are conditionally uncorrelated, this reduces to
$\sum_i p_i^2\|v_i-o\|_2^2
\operatorname{Var}[\delta s_{i,g}(b)]$.
Let $J_{>\ell}$ map the layer-$\ell$ attention output locally to final
logits and let $W_\ell^O$ be its output projection.  Replacing the Euclidean
metric in \cref{eq:output-aware-band-distortion} by
\begin{equation}
 H_\ell=(W_\ell^O)^\top J_{>\ell}^\top J_{>\ell}W_\ell^O
 \label{eq:layer-output-metric}
\end{equation}
gives a first-order logit-aware distortion.  Hence Query direction,
attention mass, Value leverage, and downstream amplification all make
$\alpha_{c,\ell,h,g}$ condition dependent.  QKSieve uses the cheaper
crossing surrogate rather than computing full downstream Jacobians.

\paragraph{Proposition (no universally optimal frozen allocation).}
There are two-band, equal-rate condition families for which every frozen
allocation fails on at least one condition, while a condition-specific
allocation succeeds on every condition at the same total rate.  One
construction lets the budget retain exactly one band: one request places
all ranking signal in band one and another places it in band two.  The same
construction can represent different layers or heads.  Keeping a fixed
Query and appending a new highest-scoring token whose distinguishing signal
lies in the other band gives a length-conditioned version.

\paragraph{Order-statistic mechanism.}
For $N$ independent uniform scores, every adjacent order-statistic spacing
has expectation
\begin{equation}
 \mathbb E[s_{(B)}-s_{(B+1)}]=\frac1{N+1}.
 \label{eq:uniform-order-gap}
\end{equation}
For a smooth score CDF $F$, with
$u_N=F^{-1}(1-B/N)$ and density $f(u_N)>0$, the probability integral
transform gives the local approximation
\begin{equation}
 m_{B,N}\asymp\frac1{Nf(u_N)}.
 \label{eq:general-order-gap}
\end{equation}
This is a mechanism, not an i.i.d.\ model of real attention.  It shows why
more competitors can tighten ranking precision even when the oracle
attention budget remains adequate; empirical margins must still be reported
per layer and head.

\paragraph{Proposition 4 (Query-distribution shift).}
Let $C'_{q,\mathrm{cal}}$ be the projected Query moment used by the
allocator, and let $C'_{q,\mathrm{dec}}$ be the moment of later decode
Queries.  For any fixed Key-error moment $C_e\succeq0$,
\begin{equation}
 \left|\tr(C'_{q,\mathrm{dec}}C_e)
       -\tr(C'_{q,\mathrm{cal}}C_e)\right|
 \le
 \|C'_{q,\mathrm{dec}}-C'_{q,\mathrm{cal}}\|_{\op}\tr(C_e).
 \label{eq:query-drift-bound}
\end{equation}
Thus prompt-tail calibration remains reliable when its projected second
moment is close to that of generation Queries.  This is a measurable
condition, not an assumption hidden inside the allocator: we report the
left-hand distortion and the covariance-drift factor separately by layer,
head, and generation position.

\paragraph{Finite-sample calibration versus genuine drift.}
Let $C'_{q,\mathrm{cal}}$ denote the population calibration moment and
$\widehat C'_{q,m}$ the moment of $m$ independent calibration Queries.
For $X_j=q'_j{q'_j}^\top-C'_{q,\mathrm{cal}}$, assume
$\|X_j\|_{\op}\le L$ almost surely and
$\|\mathbb E X_j^2\|_{\op}\le v$.
Matrix Bernstein implies that, with probability at least $1-\delta$,
\begin{equation}
 \|\widehat C'_{q,m}-C'_{q,\mathrm{dec}}\|_{\op}
 \le
 \underbrace{\sqrt{\frac{2v\log(2d/\delta)}{m}}+
 \frac{2L\log(2d/\delta)}{3m}}_{\text{finite-sample term}}
 +\underbrace{\|C'_{q,\mathrm{cal}}-C'_{q,\mathrm{dec}}\|_{\op}}
 _{\text{distribution-shift term}}.
 \label{eq:finite-query-bound}
\end{equation}
Combining this with \cref{eq:query-drift-bound} gives a high-probability
bound on allocator distortion.  Autoregressive Queries are not generally
i.i.d.; therefore \cref{eq:finite-query-bound} motivates, but does not
replace, the generation-position drift measurements in our evaluation.

\paragraph{Theorem (uniform finite-candidate allocation selection).}
Fix the QK transform and deterministic per-band quantizers before drawing
the allocation-calibration samples.  Let $\mathcal B_R$ contain the $M$
feasible assignments under \cref{eq:qmse-allocation}; for our eight bands,
five bit levels, and normalized rate at most 15, exact enumeration gives
$M=13{,}817$.  Define the separable per-pair loss
\begin{equation}
 \ell_b(q',k')=
 \sum_g\left({q'_g}^{\top}e_g(k';b_g)\right)^2
\end{equation}
and suppose $0\le\ell_b(q',k')\le L_{\rm loss}$ uniformly.
Draw $m_q$ i.i.d.\ Queries and $m_k$ i.i.d.\ Keys independently of each
other.  Let $J(b)=\mathbb E[\ell_b(q',k')]$ and let $\widehat J(b)$ be the
Cartesian empirical average over the two sample sets.
Then, with probability at least $1-\delta$,
\begin{equation}
 \sup_{b\in\mathcal B_R}|\widehat J(b)-J(b)|
 \le
 \varepsilon_{\rm sel}:=
 L_{\rm loss}\!\left[
 \sqrt{\frac{\log(4M/\delta)}{2m_q}}+
 \sqrt{\frac{\log(4M/\delta)}{2m_k}}
 \right].
 \label{eq:allocation-uniform-bound}
\end{equation}
Consequently, if $\widehat b$ minimizes $\widehat J$, then
\begin{equation}
 J(\widehat b)-\min_{b\in\mathcal B_R}J(b)
 \le2\varepsilon_{\rm sel}.
 \label{eq:allocation-selection-regret}
\end{equation}
For the full score MSE
$J_{\rm full}(b)=J(b)+\Gamma(b)$ and its minimizer
$b_{\rm full}^\star$,
\begin{equation}
 J_{\rm full}(\widehat b)-J_{\rm full}(b_{\rm full}^\star)
 \le
 2\varepsilon_{\rm sel}
 +|\Gamma(\widehat b)|+|\Gamma(b_{\rm full}^\star)|.
 \label{eq:allocation-full-regret}
\end{equation}
The Cartesian table contains $m_qm_k$ entries, but they are not independent;
the bound therefore has separate $m_q^{-1/2}$ and $m_k^{-1/2}$ terms rather
than an artificial $(m_qm_k)^{-1/2}$ rate.  The finite class enters only
through $\log M$.

This theorem applies directly to a split-sample audit in which the transform
is frozen before independent allocation samples are drawn.  Production
reuses prompt samples to estimate both the transform and the allocation, so
the theorem is not presented as an unconditional guarantee for that path.
We instead report independent calibration--held-out gaps together with the
moment and subspace perturbations above.  The bounded-loss result is also
intentionally conservative; its role is to expose the correct sources of
selection error.

\paragraph{Theorem (stability of the estimated QK-sensitive subspace).}
The preceding bound treats the QK-balanced transform as fixed.  We now
separate the additional error caused by estimating the construction moments.
Let $C_q,C_k,\widehat C_q,\widehat C_k\succ0$ be floored population and
empirical construction moments, and define
\begin{equation}
 M=C_q^{1/2}C_k^{1/2},\qquad
 \widehat M=\widehat C_q^{1/2}\widehat C_k^{1/2}.
\end{equation}
Write $\varepsilon_q=\|\widehat C_q-C_q\|_{\op}$ and
$\varepsilon_k=\|\widehat C_k-C_k\|_{\op}$, and set
\begin{align}
 a_q&=
 \frac{\varepsilon_q}
 {\sqrt{\lambda_{\min}(C_q)}
  +\sqrt{\lambda_{\min}(\widehat C_q)}},\\
 a_k&=
 \frac{\varepsilon_k}
 {\sqrt{\lambda_{\min}(C_k)}
  +\sqrt{\lambda_{\min}(\widehat C_k)}},\\
 \varepsilon_M&=
 a_q\sqrt{\|\widehat C_k\|_{\op}}
 +\sqrt{\|C_q\|_{\op}}\,a_k .
 \label{eq:qk-moment-product-perturbation}
\end{align}
Then
\begin{equation}
 \|\widehat M-M\|_{\op}\le\varepsilon_M.
 \label{eq:qk-product-stability}
\end{equation}
Let $U_r,V_r$ and $\widehat U_r,\widehat V_r$ be the leading left and
right rank-$r$ singular subspaces of $M$ and $\widehat M$, respectively.
If the population boundary gap
$\gamma_r=\sigma_r(M)-\sigma_{r+1}(M)$ satisfies
$\varepsilon_M<\gamma_r/2$, then
\begin{equation}
 \max\!\left\{
 \|\sin\Theta(\widehat U_r,U_r)\|_{\op},
 \|\sin\Theta(\widehat V_r,V_r)\|_{\op}
 \right\}
 \le \frac{2\sqrt{2}\,\varepsilon_M}{\gamma_r}.
 \label{eq:qk-subspace-stability}
\end{equation}
The floor and Query shrinkage keep the denominators in
\cref{eq:qk-moment-product-perturbation} nonzero; the singular gap determines
whether a 16-dimensional band boundary is statistically identifiable.
This result does not assert that eight prompt-tail positions always suffice.
It yields auditable conditions: report covariance estimation error,
the gaps at $r\in\{16,32,\ldots,112\}$, and the corresponding subspace
angles as the Query sample count changes.  Importantly, an inaccurate
estimate changes compression coordinates but still cannot change the
unquantized dot product, because Proposition~1 holds for any positive
definite pair of estimated moments.

\paragraph{Key reconstruction can choose the wrong bits.}
Key-only quantization minimizes
$J_K=\mathbb E\|e\|_2^2=\tr(C_e)$, whereas attention-logit distortion is
$J_{QK}=\tr(C'_qC_e)$.
These objectives are proportional for all possible error covariances only
when $C'_q$ is proportional to the identity.
When $C'_q$ is anisotropic, there exist two quantization errors for which
$J_K$ prefers the first and $J_{QK}$ prefers the second
(\cref{app:proofs}).
This is the precise reason that low Key reconstruction error is neither
necessary nor sufficient for low attention-score error.

\paragraph{Proposition 5 (cost of omitting Query covariance).}
Fix a coordinate system and feasible family of Key-error moments
$C_e(b)\succeq0$.  Suppose
$\mu I\preceq C'_q\preceq LI$ with $\mu>0$.
Let $b_K$ minimize $\tr(C_e(b))$ and let $b_{QK}$ minimize
$\tr(C'_qC_e(b))$ over the same physical rate constraint.  Then
\begin{equation}
 \tr(C'_qC_e(b_K))
 \le \kappa(C'_q)\,
      \tr(C'_qC_e(b_{QK})),
 \qquad \kappa(C'_q)=L/\mu.
 \label{eq:key-only-condition-bound}
\end{equation}
The factor becomes one for isotropic Queries and can be arbitrarily large
for anisotropic or nearly rank-deficient Queries.  This is an approximation
guarantee, not a claim that the worst case is typical; our
without-Query-covariance ablation measures the realized gap under identical
basis, quantizers, rate, and active KV.

\paragraph{Lemma (random rotation has no preferred band in expectation).}
Let $R$ be Haar-distributed on the orthogonal group and let $E_g$ select a
fixed $d_g$-dimensional coordinate band.  For any positive-semidefinite
second moment $C$,
\begin{equation}
 \mathbb E_R\tr(CRE_gR^\top)=\frac{d_g}{d}\tr(C).
 \label{eq:random-band-exchangeability}
\end{equation}
Thus a data-independent random rotation preserves exact full-precision dot
products when applied to Query and Key, but its band labels are exchangeable:
in expectation it cannot place Query--Key-sensitive directions in earlier
or higher-rate bands.  The random-rotation ablation separates the benefit of
an orthogonal change of coordinates from the benefit of QK alignment.

\paragraph{Query quantization.}
The allocator above models Key error.
At runtime the transformed Query is also quantized:
$\widehat{q'}=q'+e_q$ and $\widehat{k'}=k'-e_k$ under our sign convention.
The complete proxy error is
\begin{equation}
 {q'}^\top k'-{\widehat{q'}}^\top\widehat{k'}
 ={q'}^\top e_k-e_q^\top k'+e_q^\top e_k.
 \label{eq:query-key-quant-error}
\end{equation}
The additional term is bounded by
$\|e_q\|_2\|k'-e_k\|_2$.
INT8 Query error is reported independently; it is not silently attributed to
the mixed-bit allocator.

\subsection{From score error to attention-output error}

Let exact and proxy scores be $s_i$ and $\widehat s_i$, and define
$\delta_i=s_i-\widehat s_i$ and
$R_\delta=\max_i\delta_i-\min_i\delta_i$.
\paragraph{Lemma (softmax stability modulo a common shift).}
For $p=\softmax(s)$ and $\widehat p=\softmax(\widehat s)$,
\begin{equation}
 \mathrm{KL}(p\|\widehat p)\le\frac{R_\delta^2}{8},
 \qquad
 \|p-\widehat p\|_1\le
 \min\!\left\{2,\frac{R_\delta}{2}\right\}.
 \label{eq:softmax-range-bound}
\end{equation}
The bound depends only on the range of the score error because softmax is
invariant to a common shift.  QKSieve does not use $\widehat p$ for the final
Value aggregation, so this lemma quantifies proxy-score fidelity but does
not replace the selected-mass guarantee below.

For an exact top-$B$ token $i$ and any token $j$ outside exact top-$B$,
\begin{equation}
 s_i-s_j>R_\delta\quad\Longrightarrow\quad
 \widehat s_i>\widehat s_j.
 \label{eq:ranking-margin}
\end{equation}
Hence all exact top-$B$ tokens separated from the outside boundary by more
than $R_\delta$ are guaranteed to survive proxy top-$B$.
Unlike a raw top-$B$ overlap, this condition distinguishes high-margin core
tokens from near-tied boundary swaps.

The range can be certified directly from centered score RMSE.
Let $\bar\delta=N^{-1}\sum_i\delta_i$ and
$\eta^2=N^{-1}\sum_i(\delta_i-\bar\delta)^2$.
Then $R_\delta\le2\sqrt{N}\eta$.
If $s_{(1)}\ge\cdots\ge s_{(N)}$ and, for some $r\le B$,
\begin{equation}
 s_{(r)}-s_{(B+1)}>2\sqrt{N}\eta,
 \label{eq:rmse-core-margin}
\end{equation}
all exact top-$r$ tokens must occur in proxy top-$B$.
Although this worst-case certificate is conservative, it connects the
allocator's score-error quantity to a query-specific preserved core without
assuming Gaussian errors.

\paragraph{A less conservative count certificate.}
Let $m_r=s_{(r)}-s_{(B+1)}>0$ and
\begin{equation}
 c_r=\left\lfloor\frac{4N\eta^2}{m_r^2}\right\rfloor,
 \qquad b_r=\min\{r,c_r\}.
 \label{eq:rmse-bad-count}
\end{equation}
Even when \cref{eq:rmse-core-margin} fails, proxy top-$B$ contains at least
$r-b_r$ exact top-$r$ tokens.  Moreover, its exact attention mass is at
least
\begin{equation}
 \sum_{i=b_r+1}^{r}p_{(i)}.
 \label{eq:rmse-mass-count}
\end{equation}
The result is deterministic: at most $c_r$ tokens can have centered score
error at least $m_r/2$ in magnitude, and all remaining core/outside pairs
retain their order.  It converts average score error into a graded guarantee
rather than requiring a zero-error worst case.

\paragraph{Corollary (length amplification at fixed rate).}
Write $\eta_N$ and $m_{r,N}$ for the centered proxy-score RMSE and exact
top-$r$-to-outside margin at history length $N$.  The two certificates above
obey
\begin{equation}
 R_{\delta,N}\le2\sqrt{N}\eta_N,\qquad
 c_{r,N}\le
 \left\lfloor
 \frac{4N\eta_N^2}{m_{r,N}^2}
 \right\rfloor .
 \label{eq:length-amplification}
\end{equation}
Consequently, if $\eta_N$ and $m_{r,N}$ remain fixed while $N$ doubles, the
uniform range bound grows by $\sqrt{2}$ and the bad-count bound doubles.
Keeping either sufficient certificate constant requires, respectively,
$\eta_N=O(N^{-1/2})$ or
$\eta_N/m_{r,N}=O(N^{-1/2})$.  These are worst-case sufficient conditions,
not a claim that typical errors are independent or that quality must decay
at this rate.  They explain how an exact top-$B$ oracle can remain adequate
while a fixed-bit selector crosses a length-dependent ranking boundary.

\paragraph{Theorem (proxy top-$B$ mass inflation).}
Let $T$ be exact top-$B$, let $S$ be proxy top-$B$, and let
\begin{equation}
 \epsilon_T=\sum_{i\notin T}p_i,\qquad
 \epsilon_S=\sum_{i\notin S}p_i .
\end{equation}
For any deterministic tie breaking,
\begin{equation}
 \epsilon_S\le
 \min\!\left\{1,\exp(R_\delta)\epsilon_T\right\}.
 \label{eq:proxy-topb-mass-inflation}
\end{equation}
Thus proxy retrieval loses at most an $\exp(R_\delta)$ factor more exact
attention mass than oracle exact top-$B$.  Combining this with exact
restricted attention gives
\begin{equation}
 \|o-\widetilde o_S\|_2
 \le
 \min\!\left\{1,\exp(R_\delta)\epsilon_T\right\}
 \diam\{v_1,\ldots,v_N\}.
 \label{eq:proxy-topb-output-bound}
\end{equation}
Because $R_\delta\le2\sqrt N\,\eta$, the result also supplies a fully
deterministic, although conservative, score-RMSE-to-output certificate.
Unlike top-$B$ overlap, it discounts swaps among nearly tied low-mass tokens.

For the selected set $S$, let
$\epsilon=\sum_{i\notin S}p_i$ be omitted full-attention mass.
Restricting and renormalizing exact attention gives
\begin{equation}
 \|p-\widetilde p\|_1=2\epsilon.
 \label{eq:l1-mass}
\end{equation}
Writing $o_{\rm in}$ and $o_{\rm out}$ for conditional Value means inside and
outside $S$ yields the exact identity and bound
\begin{equation}
 \|o-\widetilde o\|_2
 =\epsilon\|o_{\rm out}-o_{\rm in}\|_2
 \le\epsilon\,\diam\{v_1,\ldots,v_N\}.
 \label{eq:output-mass}
\end{equation}
The theory therefore forms a conditional chain:
mixed-bit score error plus score margins determines a preserved core;
the actually omitted probability mass then controls the exact sparse
attention output.
It does not guarantee quality for adversarial near ties or unbounded Values.
In particular, if \cref{eq:rmse-core-margin} holds, then
\begin{equation}
 \epsilon\le1-\sum_{i=1}^{r}p_{(i)},
 \label{eq:certified-mass}
\end{equation}
which turns the score certificate into an explicit attention-output bound.

\paragraph{Proposition (shrunken Value-tail error and optimum).}
Let $Z_S,Z_T>0$ be exact selected and omitted partitions, let
$p=Z_S/(Z_S+Z_T)$, and let $\mu_S,\mu_T$ be their exact conditional Value
means.  For any tail partition estimate $\widehat Z_T\ge0$, conditional mean
estimate $\widehat\mu_T$, and $\alpha\in[0,1]$, define
\begin{equation}
 \widehat p_\alpha=
 \frac{Z_S}{Z_S+\alpha\widehat Z_T},\qquad
 \widehat o_\alpha=
 \widehat p_\alpha\mu_S+(1-\widehat p_\alpha)\widehat\mu_T.
\end{equation}
Then
\begin{equation}
 \widehat o_\alpha-o
 = (\widehat p_\alpha-p)(\mu_S-\mu_T)
 +(1-\widehat p_\alpha)(\widehat\mu_T-\mu_T),
 \label{eq:robust-tail-exact-decomposition}
\end{equation}
which immediately implies \cref{eq:tail-error-bound}.  Moreover, suppose
$\widehat Z_T=Z_T$ and
$\widehat\mu_T=\mu_T+\xi$ with $\mathbb E[\xi]=0$ and
$V=\mathbb E\|\xi\|_2^2<\infty$.  Set
$D=\|\mu_S-\mu_T\|_2^2$ and
$w_\alpha=\alpha(1-p)/(p+\alpha(1-p))$.  Then
\begin{equation}
 \mathbb E\|\widehat o_\alpha-o\|_2^2
 =((1-p)-w_\alpha)^2D+w_\alpha^2V,
\end{equation}
and the minimizers are
\begin{equation}
 w^\star=\frac{(1-p)D}{D+V},\qquad
 \alpha^\star=\frac{pD}{pD+V}.
 \label{eq:robust-tail-optimal-alpha}
\end{equation}
For $w_\alpha>0$, this correction has lower expected error than Fast
($w=0$) exactly when
\begin{equation}
 0<w_\alpha<\frac{2(1-p)D}{D+V}.
 \label{eq:robust-tail-better-than-fast}
\end{equation}

\emph{Proof.}
The full output is $o=p\mu_S+(1-p)\mu_T$.  Add and subtract
$(1-\widehat p_\alpha)\mu_T$ to obtain
\cref{eq:robust-tail-exact-decomposition}.  Under an exact tail partition,
$1-\widehat p_\alpha=w_\alpha$; expanding the squared norm and using
$\mathbb E[\xi]=0$ removes the cross term.  Differentiation with respect to
$w$ gives $w^\star$.  Solving
$w=\alpha(1-p)/(p+\alpha(1-p))$ for $\alpha$ gives
\cref{eq:robust-tail-optimal-alpha}.  Subtracting the $w=0$ error gives
$w^2(D+V)-2w(1-p)D$; its negative interval is
\cref{eq:robust-tail-better-than-fast}.  \hfill$\square$

The proposition does not claim that the proxy partition is exact or that a
single $\alpha$ is optimal for all heads.  It explains why unshrunk residual
correction can be worse when the tail estimate is noisy and why held-out
validation of the fixed $\alpha=.5$ is necessary.

\paragraph{Proposition 6 (propagation to decoder logits).}
Write the dense and sparse layer maps as $F_\ell$ and
$\widetilde F_\ell$.  Suppose on the states reached by the two executions
\begin{equation}
 \|\widetilde F_\ell(x)-\widetilde F_\ell(y)\|_2
 \le L_\ell\|x-y\|_2,\qquad
 \|\widetilde F_\ell(x)-F_\ell(x)\|_2\le a_\ell.
\end{equation}
Starting from the same hidden state, the final hidden-state error obeys
\begin{equation}
 \|\widetilde h_L-h_L\|_2
 \le\sum_{\ell=0}^{L-1}a_\ell
       \prod_{u=\ell+1}^{L-1}L_u.
 \label{eq:layer-propagation}
\end{equation}
If the local change enters only through attention, then
$a_\ell\le M_\ell
\big(\sum_h\epsilon_{\ell,h}^2\diam(V_{\ell,h})^2\big)^{1/2}$
for the downstream gain $M_\ell$ from the concatenated head outputs.
Let $\rho=\|W_U\|_{2\rightarrow\infty}$ times the right-hand side of
\cref{eq:layer-propagation}.  Then the dense and sparse next-token
distributions obey
\begin{equation}
 \mathrm{KL}(p_{\rm next}\|\widetilde p_{\rm next})\le\frac{\rho^2}{2},
 \qquad
 \|p_{\rm next}-\widetilde p_{\rm next}\|_1
 \le\min\{2,\rho\}.
 \label{eq:next-token-stability}
\end{equation}
The next-token argmax is unchanged whenever the dense logit margin exceeds
$2\rho$.
The Lipschitz product is intentionally a conditional certificate rather than
a claim that deep worst-case constants are small.

\paragraph{Lemma (tail-resolution sample complexity).}
Let $T$ contain the top $B=rN$ proxy scores and draw $m$ token positions
uniformly without replacement.  If $X$ is the number of sampled positions
in $T$, then
\begin{equation}
 m(N)=\min\!\left\{N,m_{\max},
 \max\!\left(m_0,\left\lceil\frac{c}{r(N)}\right\rceil\right)\right\},
 \quad (m_0,c,m_{\max})=(256,16,512).
 \label{eq:tail-resolution-samples}
\end{equation}
For any fixed $m$,
\begin{equation}
 \mathbb E[X]=mr,\qquad
 \Pr(X=0)=\frac{\binom{N-B}{m}}{\binom{N}{m}}
 \le (1-r)^m\le e^{-mr}.
 \label{eq:tail-sample-no-hit}
\end{equation}
Consequently, while the cap is inactive, choosing $m\ge c/r$ keeps the
expected number of target-tail anchors at least $c$ and makes the no-anchor
probability at most $e^{-c}$.  Once $m_{\max}=512$ is active, this guarantee
weakens to $mr$ anchors in expectation; the cap is a measured throughput--
accuracy tradeoff, not a length-independent theorem.  The deployment kernel uses
regular stratification with a head-dependent phase rather than an
independent random sample, so \cref{eq:tail-sample-no-hit} is the motivation
for \cref{eq:tail-resolution-samples}, not an adversarial guarantee for the
deterministic phases or the capped long-context regime.  Candidate counts and
downstream quality are therefore audited empirically.

\paragraph{Proposition (finite-sample candidate-count variation).}
The no-hit bound controls only the event that the target tail is completely
missed; it does not imply an accurate candidate count.  To expose the second
effect, assume continuous proxy scores and uniformly sampled score ranks, and
let the threshold be the $j$-th largest of $m$ samples.  The population
upper-tail mass $\widehat r$ above this order statistic obeys
\begin{equation}
 \widehat r\sim\operatorname{Beta}(j,m+1-j),\qquad
 \mathbb E[\widehat r]=\frac{j}{m+1},
 \label{eq:sample-quantile-beta}
\end{equation}
and
\begin{equation}
 \operatorname{Var}(\widehat r)
 =\frac{j(m+1-j)}{(m+1)^2(m+2)}.
 \label{eq:sample-quantile-variance}
\end{equation}
For $j\approx mr=c$ and small $r$, the relative standard deviation of the
realized candidate count $N\widehat r$ is therefore
\begin{equation}
 \frac{\sqrt{\operatorname{Var}(N\widehat r)}}
 {\mathbb E[N\widehat r]}
 =
 \sqrt{\frac{m+1-j}{j(m+2)}}
 \approx \frac{1}{\sqrt{c}}.
 \label{eq:candidate-count-relative-std}
\end{equation}
Thus $c=16$ gives approximately $25\%$ relative standard deviation even
though its no-hit probability is tiny; $c=64$ and $c=256$ reduce the
diagnostic value to approximately $12.5\%$ and $6.25\%$.  This calculation
does not assume that the deterministic stratified phases are i.i.d.; it is a
finite-sample reference model whose prediction is checked against per-head
candidate-count dispersion.  It motivates reporting both the mean and
headwise range and testing larger $c$ or a conservative overselect-and-trim
selector at very long context.

\paragraph{Headwise stragglers.}
The mean candidate count is not the relevant system statistic when a ragged
consumer sizes work from the largest layer-head.  The same order-statistic
model gives a useful tail prediction.  In the rare-tail limit
$m\rightarrow\infty$, $r\rightarrow0$, and $mr=c$, the normalized realized
tail mass satisfies
\begin{equation}
 Z=\frac{\widehat r}{r}
 \ \Longrightarrow\ \frac{Y}{c},
 \qquad Y\sim\operatorname{Gamma}(c,1).
 \label{eq:quantile-gamma-limit}
\end{equation}
Consequently, for $\epsilon>0$, a Chernoff bound and a union bound over $H$
layer-heads give
\begin{equation}
 \Pr\!\left[
   \max_{h\le H}\frac{\widehat C_h}{B}\ge 1+\epsilon
 \right]
 \ \lesssim\
 H\exp\{-c[\epsilon-\log(1+\epsilon)]\},
 \qquad \widehat C_h=N\widehat r_h .
 \label{eq:quantile-headwise-straggler}
\end{equation}
For $H=288$ and failure probability $.05$, the corresponding upper factors
$1+\epsilon$ are approximately $2.43$, $1.61$, and $1.28$ for
$c=16,64,256$.  This is an asymptotic reference bound rather than a
certificate for deterministic stratification, but it predicts the observed
system effect: increasing $c$ can reduce maximum candidate capacity, split
count, and tail latency enough to offset the extra sampled comparisons.

\subsection{Complexity and request-level crossover}

Ignoring constants, dense decode attention costs $4H_qNd$ floating-point
operations per layer for QK and AV.
\method costs
\begin{equation}
 2H_qd^2+c_{\rm idx}H_qN R+
 C_{\rm select}+4H_qB(N)d,
 \label{eq:cost}
\end{equation}
where $R=240$ packed bits per token/KV head and $c_{\rm idx}$ reflects
low-bit scan throughput.  The reference profile has
$C_{\rm select}=c_{\rm topk}H_qN$; the deployment profile replaces this by
the stratified sample and in-block order statistic,
$C_{\rm select}=O(H_qm(N)R+H_qm(N)\log^2m(N))$, followed by candidate writes
during the mandatory packed scan.
The scan is still $\Theta(N)$; the gain comes from reading and processing a
small packed Key index and capping exact attention, not from sublinear search.

Let $I$ be visible one-time index construction cost and let
$t_{\rm full}$ and $t_{\method}$ be synchronized steady decode times.
For $G$ generated tokens,
\begin{equation}
 T_{\rm full}=P_{\rm full}+Gt_{\rm full},\qquad
 T_{\method}=P_{\rm full}+I+Gt_{\method}.
 \label{eq:request-cost}
\end{equation}
When $t_{\method}<t_{\rm full}$, the break-even output length is
\begin{equation}
 G^\star=\frac{I}{t_{\rm full}-t_{\method}}.
 \label{eq:break-even}
\end{equation}
This accounting explains why a strong 128K decode speedup can coexist with a
smaller short-output request speedup, and why 2K--8K contexts can remain
slower until index scan and top-$k$ are sufficiently amortized.

\paragraph{Proposition 7 (exact crossover and Amdahl ceiling).}
Because history grows during generation, let $f_t$ and $s_t$ be measured
Full and \method decode costs at output step $t$, and let $I_{\rm net}$ be
index construction cost not hidden under prefill.  After $G$ tokens,
\method is faster at request level if and only if
\begin{equation}
 \sum_{t=1}^{G}(f_t-s_t)>I_{\rm net}.
 \label{eq:exact-break-even}
\end{equation}
Thus if $s_t\ge f_t$ at every step, no output length can amortize the method;
using a Full fallback would hide rather than solve this regime.
For one decode step, write
$f=t_{\rm base}+t_{\rm attn}$ and
$s=t_{\rm base}+\widetilde t_{\rm attn}$, with all retrieval overhead
included in $\widetilde t_{\rm attn}$.  Then
\begin{equation}
 \mathrm{speedup}=
 \frac{t_{\rm base}+t_{\rm attn}}
      {t_{\rm base}+\widetilde t_{\rm attn}}
 \le
 \frac{t_{\rm base}+t_{\rm attn}}{t_{\rm base}}
 =\frac{1}{1-\phi_{\rm attn}},
 \label{eq:amdahl-ceiling}
\end{equation}
where $\phi_{\rm attn}=t_{\rm attn}/f$.
The last expression is the zero-attention ceiling.  The measured decomposition
in \cref{fig:speed,tab:deployment-breakdown} shows why it is not an attainable
kernel claim.

\paragraph{Proposition 8 (fused Query-preparation certificate).}
Let $\widehat z$ and $\widetilde z$ be the dequantized transformed Queries
produced by the reference projection--quantization sequence and by a fused
kernel, respectively.  For fixed decoded proxy Keys $\widehat k_i$, define
\begin{equation}
 \epsilon_{\rm fuse}
 =
 \|\widetilde z-\widehat z\|_2
 \max_i\|\widehat k_i\|_2 .
 \label{eq:fused-query-epsilon}
\end{equation}
Then every proxy score changes by at most $\epsilon_{\rm fuse}$.
Consequently, the top-$B$ set is invariant whenever the reference boundary
gap satisfies
\begin{equation}
 \widehat s_{(B)}-\widehat s_{(B+1)}
 >2\epsilon_{\rm fuse}.
 \label{eq:fused-query-margin}
\end{equation}
If the fused kernel produces exactly the same INT8 codes and scales as the
reference sequence, proxy scores and the top-$B$ set are identical, modulo
ordering among exact ties.  Given the same selected set and original exact
K/V, the two sparse-attention outputs are identical in real arithmetic.
Finite-precision GPU reductions can still differ slightly, so kernel
promotion additionally requires measured code/scale, score, selected-set,
and final-output agreement rather than projection latency alone.
